% !PRINTED 504
\weberpage{504}
%<<CONT
nochmals durch $F(t)$, d. h. $N(t-\omega)$ ist durch $F(t)^2$ theilbar.
Durch Fortsetzung dieses Schlussverfahrens erkennt man, dass
$N(t-\omega)$ eine Potenz von $F(t)$ sein muss, und dass folglich der
Grad von $F(t)$ ein Theiler von $n$ ist (Bd. I, \S.\,144).

Wir fassen dies noch in dem Satz zusammen:

8. \gesperrt{Jedes Functional in $\Omega$ ist die Wurzel einer und
nur einer irreduciblen Gleichung $F(t) = 0$ in $\overline{R}$
und $N(t-\omega)$ ist eine Potenz von $F(t)$.}

\begin{center}
\S.\ 137.

\gesperrt{Ganze Functionale.}
\end{center}

1. \textbf{Definition:} \gesperrt{Ein rationales Functional soll ganz
genannt werden, wenn sein absoluter Werth eine ganze Zahl ist.}

Es ergiebt sich aus dieser Definition zunächst, dass die
Summe, die Differenz und das Product von zwei und folglich
auch von beliebig vielen ganzen rationalen Functionalen wieder
ganz sind. Für das Product ist dies eine unmittelbare Folge
des Satzes \S.\,136, 3. Um aber den Beweis für die Summe und
die Differenz zu führen, stellen wir zwei ganze rationale Func\-
tionale $A_1, A_2$ so durch gebrochene Functionen dar, dass sie eine
primitive Form als gemeinschaftlichen Nenner erhalten:
\[
A_1 = a_1\,\frac{E_1}{E}, \quad A_2 = a_2\,\frac{E_2}{E}.
\]
Dann ist
\[
A_1 \pm A_2 = \frac{a_1 E_1 \pm a_2 E_2}{E},
\]
da nun $a_1, a_2$ ganze Zahlen sind, so ist der Theiler der ganzen
Function $a_1 E_1 \pm a_2 E_2$ der absolute Werth von $A_1 \pm A_2$. Dieser
ist eine ganze Zahl, also $A_1 \pm A_2$ ein ganzes Functional.

Für constante rationale Functionale, d. h. für rationale
Zahlen, giebt die Definition 1. die \gesperrt{ganzen rationalen Zahlen.}

Wir stellen ferner, ebenso wie im \S.\,133 für die Zahlen,
folgende Definition der ganzen Functionale des Körpers $\Omega$ auf.

2. \textbf{Definition:} \gesperrt{Ein Functional $\omega$ aus $\Omega$ heisst
ganz, wenn es die Wurzel einer Gleichung}
\[
\Phi(t) = t^m + A_1 t^{m-1} + A_2 t^{m-2} + \cdots + A_m = 0
\]
%CONT>>
